\section{Exponential bounds on invariant subspaces}

The inverse metric comparison in \cite[Section~8]{OpenAI2026AreaLaw}
uses exponential inequalities on invariant subspaces. All matrices below
act on finite-dimensional complex inner product spaces, which may have
dimension zero. Matrix order means positive semidefinite order.

\begin{theorem}[Exponential order for commuting matrices]
  \label{thm:commuting_exponential_order}
  \lean{Matrix.IsHermitian.exp_le_exp_of_commute}
  Let $A,B$ be commuting Hermitian matrices. If $A\leq B$, then
  $e^A\leq e^B$.
\end{theorem}
\begin{proof}
  Set $D=B-A\geq0$. Scalar functional calculus gives $e^D\geq I$.
  Since $A$ and $D$ commute, $e^B-e^A=e^A(e^D-I)\geq0$.
\end{proof}

\begin{theorem}[Functional calculus on an invariant range]
  \label{thm:invariant_range_functional_bound}
  \lean{Matrix.IsHermitian.compression_cfc_le_of_lower_bound}
  Let $A=A^*$ and let $P=P^*=P^2$ commute with $A$. Suppose that
  $bP\leq PAP$ for some $b\in\mathbb R$. If $f\colon\mathbb R\to\mathbb R$
  and $c\in\mathbb R$ satisfy $f(x)\leq c$ whenever $x\geq b$, then
  $Pf(A)P\leq cP$.
\end{theorem}
\begin{proof}
  The Hermitian matrix $B=PA+b(I-P)$ satisfies $B\geq bI$.
  Functional calculus on its finite spectrum gives $f(B)\leq cI$.
  Since $PA=PB$, the restrictions of $f(A)$ and $f(B)$ to the range
  of $P$ agree. Compressing the inequality gives the result.
\end{proof}

\begin{corollary}[Negative exponential on an invariant range]
  \label{cor:invariant_range_negative_exp}
  \lean{Matrix.IsHermitian.compression_exp_neg_smul_le_of_lower_bound}
  Under the hypotheses on $A,P,b$ in
  Theorem~\ref{thm:invariant_range_functional_bound}, every $a\geq0$
  satisfies $Pe^{-aA}P\leq e^{-ab}P$.
\end{corollary}
\begin{proof}
  \uses{thm:invariant_range_functional_bound}
  Apply the theorem to $f(x)=e^{-ax}$, which is at most $e^{-ab}$
  on $[b,\infty)$.
\end{proof}

\begin{theorem}[Exponential comparison on fixed vectors]
  \label{thm:fixed_vector_exp_comparison}
  \lean{Matrix.IsHermitian.re_dotProduct_exp_add_le_of_compressed_exp_le_smul}
  Let $A$ and $P$ be Hermitian, let $B$ be a complex matrix, and assume
  $AP=PA$ and $AB=BA$. If $Pe^BP\leq cP$ for $c\in\mathbb R$, then
  every vector $v$ with $Pv=v$ satisfies
  $\operatorname{Re}(v^*e^{A+B}v)\leq c\operatorname{Re}(v^*e^Av)$.
\end{theorem}
\begin{proof}
  \uses{thm:hermitian_compression_quadratic}
  Set $C=e^{A/2}$. Commutation gives $P(Cv)=Cv$, and the fixed-vector
  quadratic bound applied to $Cv$ gives
  $\operatorname{Re}((Cv)^*e^B(Cv))\leq c\|Cv\|^2$.
  Since $C=C^*$, $Ce^BC=e^{A+B}$ and $C^2=e^A$, this is the desired
  inequality.
\end{proof}

\begin{corollary}[Removing a compressed exponential contraction]
  \label{cor:fixed_vector_exp_contraction}
  \lean{Matrix.IsHermitian.re_dotProduct_exp_add_le_of_compressed_exp_le}
  Under the matrix hypotheses of Theorem~\ref{thm:fixed_vector_exp_comparison},
  $Pe^BP\leq P$ implies
  $\operatorname{Re}(v^*e^{A+B}v)\leq\operatorname{Re}(v^*e^Av)$
  for every $Pv=v$.
\end{corollary}
\begin{proof}
  \uses{thm:fixed_vector_exp_comparison}
  Take $c=1$.
\end{proof}

\begin{theorem}[Compressed exponential comparison]
  \label{thm:compressed_exp_comparison}
  \lean{Matrix.IsHermitian.compression_exp_add_le_of_compressed_exp_le}
  Let $A,B$ be Hermitian, let $P=P^*=P^2$, and suppose that $AP=PA$
  and $AB=BA$. If $Pe^BP\leq cP$ for $c\in\mathbb R$, then
  $Pe^{A+B}P\leq cPe^AP$.
\end{theorem}
\begin{proof}
  \uses{thm:fixed_vector_exp_comparison}
  Apply the fixed-vector comparison to $Pv$ for every $v$.
  The two compressed matrices are Hermitian, so their quadratic-form
  comparison is equivalent to their matrix order.
\end{proof}

\begin{theorem}[Exponential comparison for a supported density]
  \label{thm:supported_density_exp_comparison}
  \lean{Matrix.IsHermitian.re_trace_mul_exp_add_le_of_compressed_exp_le}
  Under the hypotheses of Theorem~\ref{thm:compressed_exp_comparison},
  let $\rho\geq0$ satisfy $P\rho=\rho$. Then
  $\operatorname{Re}\operatorname{tr}(\rho e^{A+B})
  \leq c\operatorname{Re}\operatorname{tr}(\rho e^A)$.
  No trace normalization is required.
\end{theorem}
\begin{proof}
  \uses{thm:compressed_exp_comparison}
  Taking adjoints gives $\rho P=\rho$. Pair the compressed matrix
  inequality with $\rho$, using positivity of the trace pairing.
  Cyclicity and the two support identities remove both projections.
\end{proof}
